% Generated draft: requires mathematical and source review.
\documentclass[11pt]{article}
\usepackage[margin=25mm]{geometry}
\usepackage{fontspec}
\setmainfont{Latin Modern Roman}
\usepackage{amsmath,amssymb,mathtools,unicode-math}
\setmathfont{Latin Modern Math}
\usepackage{enumitem,needspace,xurl,hyperref}
\setlength{\parindent}{0pt}
\setlength{\parskip}{5pt}
\setlength{\emergencystretch}{3em}
\newcommand{\N}{\mathbb{N}}
\newcommand{\Z}{\mathbb{Z}}
\newcommand{\Q}{\mathbb{Q}}
\newcommand{\R}{\mathbb{R}}
\newcommand{\C}{\mathbb{C}}
\newcommand{\F}{\mathbb{F}}
\newcommand{\A}{\mathbb{A}}
\newcommand{\PP}{\mathbb{P}}
\newcommand{\E}{\mathbb{E}}
\newcommand{\eps}{\varepsilon}
\newcommand{\dd}{\,\mathrm{d}}
\begin{document}
\section*{UnsolvedMath ID: 701665}
\textbf{Source problem code:} catalogue entry
\subsection{Definitions and mathematical statement}
A group \(G\) is finitely generated if it has a finite generating set \(X\). Its word problem is solvable if there is an algorithm which, for every finite word in the alphabet \(X\cup X^{-1}\), terminates and decides whether that word represents the identity of \(G\). This property is independent of the choice of finite generating set. A group is finitely presented if it admits a presentation \(\langle Y\mid R\rangle\) with both \(Y\) and \(R\) finite. A group \(S\) is simple if \(S\ne\{1\}\) and its only normal subgroups are \(\{1\}\) and \(S\). An embedding is an injective group homomorphism.

\textbf{Question (W. Boone and G. Higman; Kourovka Notebook Problem 20.7, proposed by J. Belk).} For every finitely generated group \(G\) with solvable word problem, does there exist a finitely presented simple group \(S\) and an injective homomorphism
\[
\iota:G\hookrightarrow S?
\]
There is no requirement here that \(G\) or \(S\) be infinite.

\textbf{Supplied background.} The source states that every such \(G\) embeds into a simple subgroup \(H\) of a finitely presented group \(P\), so that \(G\hookrightarrow H\leq P\). That background assertion does not require \(H\) to be finitely presented or \(P\) to be simple. The registry labels the question open; this label and the background assertion are supplied source information, not an independently verified literature assessment.
\subsection{Short English statement}
Does every finitely generated group whose word problem can be decided by an algorithm embed into a group that is both finitely presented and simple? The supplied background gives an embedding into a simple subgroup of a finitely presented group, which does not itself provide the requested target.
\subsection{Sources}
\begin{itemize}
\item Local UnsolvedMath registry, permanent ID \(701665\), exact source code \texttt{catalogue entry}, title \emph{Kourovka Notebook Problem 20.7}. The supplied record labels the problem open and records J. Belk as proposer, with proposed year 2022.
\item \emph{The Kourovka Notebook (Unsolved Problems in Group Theory)}, Problem 20.7: 20th issue (2022) and the supplied 21st issue edition (2026), arXiv:1401.0300v48. The collection description names E. I. Khukhro and V. D. Mazurov as editors. The question is attributed in the statement to W. Boone and G. Higman.
\item W. Boone and G. Higman, \emph{J. Austral. Math. Soc.} \textbf{18}, no. 1 (1974), 41--53. This is the supplied reference for the assertion that every finitely generated group with solvable word problem embeds into a simple subgroup of a finitely presented group; no article title is supplied.
\end{itemize}
\end{document}
